\chapter{Computation of volume}

\section{The Lobachevsky function}

\begin{definition}[Lobachevsky function]
\label{thurston-ch07-lobachevsky-function}
\group{lobachevsky}
\source{thurston:page-0171,thurston:page-0172}
\notready
Define
\[
\Lambda(\theta)=-\int_0^\theta\log|2\sin u|\,du.
\]
Its derivatives are $\Lambda'(\theta)=-\log|2\sin\theta|$ and
$\Lambda''(\theta)=-\cot\theta$ wherever defined.  For $0\leq\theta\leq\pi$,
the dilogarithm $\psi(z)=\sum_{n\geq1}z^n/n^2$ satisfies
\[
\psi(e^{2i\theta})-\psi(1)=-\theta(\pi-\theta)+2i\Lambda(\theta).
\]
\end{definition}

\begin{lemma}[Fourier expansion of the Lobachevsky function]
\label{thurston-ch07-lobachevsky-fourier}
\group{lobachevsky}
\source{thurston:page-0172}
\notready
The Lobachevsky function has the uniformly convergent expansion
\[
\Lambda(\theta)=\frac12\sum_{n=1}^{\infty}\frac{\sin(2n\theta)}{n^2}.
\]
\end{lemma}
\begin{proof}
\sketch
Take imaginary parts in the dilogarithm identity of the preceding definition.  Uniform convergence extends the identity to the endpoints.
\end{proof}

\begin{lemma}[Periodicity and oddness of the Lobachevsky function]
\label{thurston-ch07-lobachevsky-symmetries}
\group{lobachevsky}
\source{thurston:page-0172}
\notready
The function $\Lambda$ is odd and has period $\pi$:
$\Lambda(-\theta)=-\Lambda(\theta)$ and $\Lambda(\theta+\pi)=\Lambda(\theta)$.
\end{lemma}
\begin{proof}
\sketch
The derivative is $\pi$-periodic, while the endpoint values forced by the Fourier series are $\Lambda(0)=\Lambda(\pi)=0$; integration gives both assertions.
\end{proof}

\begin{lemma}[Multiplication identity]
\label{thurston-ch07-lobachevsky-multiplication}
\group{lobachevsky}
\source{thurston:page-0172,thurston:page-0173}
\notready
For every nonzero integer $n$,
\[
\Lambda(n\theta)=n\sum_{j\bmod |n|}\Lambda(\theta+j\pi/n).
\]
In particular, $\frac12\Lambda(2\theta)=\Lambda(\theta)+\Lambda(\theta+\pi/2)$ and $\frac32\Lambda(\pi/3)=\Lambda(\pi/6)$.
\end{lemma}
\begin{proof}
\sketch
Factor $2\sin(nu)$ into the shifted factors $2\sin(u+j\pi/n)$, integrate logarithms, and use oddness and periodicity to extend from positive $n$ to all nonzero integers.
\end{proof}

\begin{conjecture}[Rational relations among Lobachevsky values]
\label{thurston-ch07-lobachevsky-relations-conjecture}
\group{lobachevsky}
\source{thurston:page-0173}
\notready
For rational multiples of $\pi$, every rational linear relation among values of $\Lambda$ follows from oddness, period $\pi$, and the multiplication identity.  Equivalently, for $N\geq3$, the values
\[
\{\Lambda(\pi j/N):0<j<N/2,\ \gcd(j,N)=1\}
\]
are rationally linearly independent, so their span has dimension $\varphi(N)/2$.
\end{conjecture}

\begin{proposition}[Bernoulli expansion]
\label{thurston-ch07-lobachevsky-bernoulli-expansion}
\group{lobachevsky}
\source{thurston:page-0174}
\notready
Let $B_n$ denote the Bernoulli numbers in the standard convention, so that $B_2=1/6$ and $B_4=-1/30$.  For $0<|\theta|\leq\pi$,
\[
\Lambda(\theta)=\theta\left(1-\log|2\theta|+\sum_{n=1}^{\infty}\frac{(-1)^{n+1}B_{2n}}{2n}\frac{(2\theta)^{2n}}{(2n+1)!}\right).
\]
The series converges absolutely throughout $|\theta|\leq\pi$, and its right-hand side extends continuously to $0$ with value $\Lambda(0)=0$.
\end{proposition}

\section{Ideal simplices and cones}

\begin{definition}[Ideal simplex]
\label{thurston-ch07-ideal-simplex}
\group{ideal-simplex-volume}
\source{thurston:page-0174}
\notready
An ideal simplex $\Sigma_{\alpha,\beta,\gamma}$ is a hyperbolic tetrahedron whose vertices lie at infinity.  Its opposite dihedral angles agree, and the three angles satisfy $\alpha+\beta+\gamma=\pi$.
\end{definition}

\begin{lemma}[Volume of a right-angled auxiliary simplex]
\label{thurston-ch07-right-angled-simplex-volume}
\group{ideal-simplex-volume}
\source{thurston:page-0174,thurston:page-0175,thurston:page-0176}
\notready
Let $S_{\alpha,\pi/2-\alpha,\gamma}$ have three right dihedral angles, one ideal vertex, and remaining angles $\alpha,\pi/2-\alpha,\gamma$. Then
\[
\operatorname{Vol}(S_{\alpha,\pi/2-\alpha,\gamma})=
\frac14\bigl[\Lambda(\alpha+\gamma)+\Lambda(\alpha-\gamma)+2\Lambda(\pi/2-\alpha)\bigr].
\]
When $\gamma=\alpha$, this becomes $\operatorname{Vol}(S_{\alpha,\pi/2-\alpha,\alpha})=\frac12\Lambda(\alpha)$.
\end{lemma}
\begin{proof}
\sketch
In the upper half-space model place the ideal vertex at infinity and the base on the unit hemisphere.  Projecting to the horizontal plane gives $0\leq x\leq\cos\gamma$, $0\leq y\leq x\tan\alpha$.  Integrating $z^{-3}$ above the hemisphere gives $\int_T dx\,dy/[2(1-x^2-y^2)]$; integrate in $y$, then substitute $x=\cos\theta$ and identify the resulting logarithmic integral with four Lobachevsky values.  The special case follows from the symmetry and multiplication identities.
\end{proof}

\begin{theorem}[Volume of an ideal tetrahedron]
\label{thurston-ch07-ideal-tetrahedron-volume}
\group{ideal-simplex-volume}
\source{thurston:page-0174,thurston:page-0176,thurston:page-0177}
\notready
\dcref{Theorem 7.2.1}
\[
\operatorname{Vol}(\Sigma_{\alpha,\beta,\gamma})=\Lambda(\alpha)+\Lambda(\beta)+\Lambda(\gamma).
\]
\end{theorem}
\begin{proof}
\sketch
Drop the perpendicular from an ideal vertex to the opposite face and join its foot to the three base edges.  The resulting six pieces have three right dihedral angles; paired pieces are of the auxiliary type, with volumes $\frac12\Lambda(\alpha)$, $\frac12\Lambda(\beta)$, and $\frac12\Lambda(\gamma)$.  Summing them proves the formula, with signed pieces covering the case where the foot lies outside the base.
\end{proof}

\begin{theorem}[Volume of an ideal polygonal cone]
\label{thurston-ch07-ideal-polygonal-cone-volume}
\group{ideal-cone-volume}
\source{thurston:page-0178,thurston:page-0179}
\notready
\dcref{Theorem 7.2.3}
For an ideal cone $\Sigma_{\alpha_1,\ldots,\alpha_n}$ on a planar $n$-gon, the base dihedral angles satisfy
\[
\sum_{i=1}^n\alpha_i=\pi,
\qquad
\operatorname{Vol}(\Sigma_{\alpha_1,\ldots,\alpha_n})=\sum_{i=1}^n\Lambda(\alpha_i).
\]
\end{theorem}
\begin{proof}
\sketch
Induct on $n$, splitting the base polygon into a triangle and an $(n-1)$-gon.  The new angles are $\beta$ and $\pi-\beta$; apply the induction hypothesis to both cones and cancel $\Lambda(\pi-\beta)=-\Lambda(\beta)$.
\end{proof}

\begin{example}[Basic link-complement volumes]
\label{thurston-ch07-basic-link-volumes}
\group{ideal-simplex-volume}
\source{thurston:page-0178,thurston:page-0179}
\notready
Two regular ideal tetrahedra give the figure-eight knot complement volume $6\Lambda(\pi/3)=2.02988\ldots$.  Two cones on regular quadrilaterals give the Whitehead link complement volume $8\Lambda(\pi/4)=3.66386\ldots$, and two regular ideal octahedra give the Borromean-rings complement volume $16\Lambda(\pi/4)=7.32772\ldots$.
\end{example}

\section{Reflection polyhedra}

\begin{definition}[Symmetric ideal prism]
\label{thurston-ch07-symmetric-ideal-prism}
\group{reflection-polyhedra}
\source{thurston:page-0179,thurston:page-0180}
\notready
Let $\mathcal N_{\alpha,\beta}$ be the polyhedron obtained by joining corresponding vertices of congruent regular ideal $n$-gons in two planes symmetrically.  Its side faces meet the bases at angle $\alpha$, meet each other at angle $\beta$, and
\[
2\alpha+\beta=\pi.
\]
\end{definition}

\begin{proposition}[Volume of a symmetric ideal prism]
\label{thurston-ch07-symmetric-ideal-prism-volume}
\group{reflection-polyhedra}
\source{thurston:page-0180,thurston:page-0181}
\notready
\[
\operatorname{Vol}(\mathcal N_{\alpha,\beta})
=n\left(2\Lambda(\beta/2)+\Lambda(\alpha+\pi/n)+\Lambda(\alpha-\pi/n)\right).
\]
\end{proposition}
\begin{proof}
\sketch
Subdivide into $n$ congruent sectors and add the cone on the upper triangular face while subtracting the equal-volume cone on the lower face.  The resulting quadrilateral cone has side angles $\beta/2$ and front/back angles $\alpha\pm\pi/n$, by the ideal polygonal-cone angle sum and volume formula.
\end{proof}

\begin{theorem}[Infinitely many reflection-quotient manifolds]
\label{thurston-ch07-reflection-quotient-infinite-family}
\group{reflection-polyhedra}
\source{thurston:page-0181}
\notready
\dcref{Theorem 7.3.1}
There are infinitely many oriented finite-volume hyperbolic three-manifolds whose volumes are finite rational sums of $\Lambda(\theta)$ with $\theta/\pi\in\mathbb Q$.
\end{theorem}
\begin{proof}
\sketch
For $\alpha,\beta$ integral submultiples of $\pi$, reflections in the prism faces generate a discrete group.  The relation $2\alpha+\beta=\pi$ leaves $(\alpha,\beta)=(\pi/3,\pi/3)$ or $(\pi/4,\pi/2)$.  Torsion-free orientation-preserving finite-index subgroups produce manifolds, and varying $n>4$ gives infinitely many cases; the volume proposition supplies the stated form.
\end{proof}

\section{Arithmetic constructions}

\begin{definition}[Bianchi group and discriminant]
\label{thurston-ch07-bianchi-group-discriminant}
\group{arithmetic-volume}
\source{thurston:page-0181,thurston:page-0182}
\notready
For square-free $d\geq1$, let $K=\mathbb Q(\sqrt{-d})$, let $\mathcal O_d$ be its ring of integers, and define
\[
D=\begin{cases}d&d\equiv3\pmod4,\\4d&\text{otherwise.}\end{cases}
\]
The group $\operatorname{PSL}(2,\mathcal O_d)$ is discrete in $\operatorname{PSL}(2,\mathbb C)\cong\operatorname{Isom}^+(H^3)$; every torsion-free finite-index subgroup gives an oriented finite-volume hyperbolic $3$-manifold.  The lattice $\mathbb C/\mathcal O_d$ has fundamental area $\sqrt D/2$.
\end{definition}

\begin{definition}[Dedekind zeta function and quadratic symbol]
\label{thurston-ch07-dedekind-zeta-quadratic-symbol}
\group{arithmetic-volume}
\source{thurston:page-0182,thurston:page-0183}
\notready
For $K$ with ring of integers $\mathcal O_K$,
\[
\zeta_K(s)=\sum_{\mathfrak a}\frac1{N(\mathfrak a)^s}=\prod_{\mathfrak p}(1-N(\mathfrak p)^{-s})^{-1}.
\]
The quadratic symbol $\left(\frac{-D}{n}\right)$ is completely multiplicative, is $0$ for primes dividing $D$, is $1$ at $n=1$, and for odd primes records whether $-D$ is a quadratic residue modulo $p$; at $p=2$ it is $+1$ for $-D\equiv1\pmod8$ and $-1$ for $-D\equiv5\pmod8$.
\end{definition}

\begin{theorem}[Bianchi-group covolume]
\label{thurston-ch07-bianchi-covolume}
\group{arithmetic-volume}
\source{thurston:page-0182}
\notready
\dcref{Theorem 7.4.1}
\[
\operatorname{Vol}(H^3/\operatorname{PSL}(2,\mathcal O_d))
=\frac{D^{3/2}}{24}\frac{\zeta_{\mathbb Q(\sqrt{-d})}(2)}{\zeta_{\mathbb Q}(2)}.
\]
\end{theorem}

\begin{theorem}[Arithmetic Lobachevsky volume formula]
\label{thurston-ch07-arithmetic-lobachevsky-volume}
\group{arithmetic-volume}
\uses{thurston-ch07-bianchi-covolume,thurston-ch07-lobachevsky-fourier}
\source{thurston:page-0182,thurston:page-0183}
\notready
With $D$ and the quadratic symbol as above,
\[
\operatorname{Vol}(H^3/\operatorname{PSL}(2,\mathcal O_d))
=\frac D{12}\sum_{k\bmod D}\left(\frac{-D}{k}\right)\Lambda(\pi k/D).
\]
\end{theorem}
\begin{proof}
\sketch
Hecke's identity expresses the zeta quotient as $\sum_{n>0}(\frac{-D}{n})n^{-s}$.  Apply the finite Fourier transform identity
\[
\sum_{k\bmod D}\left(\frac{-D}{k}\right)e^{2\pi ikn/D}=\sqrt{-D}\left(\frac{-D}{n}\right)
\]
at $s=2$, take imaginary parts using the Fourier series of $\Lambda$, and multiply by $D/24$.
\end{proof}

\begin{example}[Arithmetic link volumes]
\label{thurston-ch07-arithmetic-link-volumes}
\group{arithmetic-volume}
\uses{thurston-ch07-arithmetic-lobachevsky-volume}
\source{thurston:page-0182,thurston:page-0183}
\notready
For $d=3$, $\mathcal O_d=\mathbb Z[\omega]$ with $\omega=(-1+\sqrt{-3})/2$, and the Bianchi orbifold has volume $\frac12\Lambda(\pi/3)$.  Its index-$12$ figure-eight subgroup therefore has volume $6\Lambda(\pi/3)$; analogous index-$12$ and index-$24$ Gaussian-integer subgroups give the Whitehead and Borromean-rings volumes above.
\end{example}
